%% ============================================================
%% Section 6: Future Research Directions
%% Source material: C15 (§6)
%% ============================================================

\section{Future Research Directions}
\label{sec:future}

The preceding sections reveal a persistent gap between AI capabilities and the protective infrastructure designed to manage them. Six critical research frontiers remain largely unsolved, each characterized by a current gap, its significance, and concrete research paths forward.


%% ------------------------------------------------------------
\subsection{Securing Autonomous AI Agents}
\label{sec:future-agents}
%% ------------------------------------------------------------

\emph{Gap:} AI agents that autonomously invoke tools, access external data, and delegate tasks introduce attack surfaces that no existing framework adequately addresses. OWASP's Top~10 for LLM Applications 2025 identifies ``Excessive Agency'' as a top risk~\cite{owasp2025}, and NIST launched a dedicated AI Agent Standards Initiative in February 2026~\cite{nist2026agentstandards}, signaling recognition that current frameworks are insufficient. Deng et al.\ provide the first comprehensive threat taxonomy for agentic systems~\cite{deng2025agentthreat}.

\emph{Why most urgent:} Agents are already deployed in production generating billions in revenue, but security practices lag severely. Inter-agent trust exploitation achieves an 82.4\% attack success rate---far exceeding direct prompt injection at 41.2\%~\cite{darkside2025agents}.

\emph{Research directions:}
\begin{enumerate}[nosep,leftmargin=1.5em]
    \item \textbf{Formal agent trust models.} No mathematically complete trust calculus exists for composing trust across delegation chains. Current proposals formalize trust as an operational variable with strength, scope, and revocability dimensions~\cite{trustparadox2025}, and explore decentralized identities (DIDs) and verifiable credentials for agents~\cite{rodriguez2025did}. Dynamic trust revocation with provable guarantees under adversarial conditions remains unsolved.
    \item \textbf{Secure orchestration across agents and MCP servers.} An empirical study of 1,899 open-source MCP servers found 7.2\% harbor general vulnerabilities and 5.5\% exhibit MCP-specific tool poisoning~\cite{hasan2025mcpglance}. B\"uhler et al.\ propose AgentBox with manifest-based permissions and sandbox enforcement~\cite{buhler2025agentbox}. Critical gaps include MCP's lack of user-context propagation (enabling confused deputy attacks) and absent standardized cross-server authentication.
    \item \textbf{Runtime verification of agent behavior.} AgentSpec introduces a lightweight DSL preventing $>$90\% of unsafe code executions~\cite{wang2025agentspec}. Pro2Guard extends this with probabilistic model checking to anticipate unsafe states~\cite{pro2guard2025}. No current system achieves both trajectory-based multi-step safety analysis and real-time enforcement without prohibitive overhead.
    \item \textbf{Secure multi-agent collaboration.} Adapting sandboxing, information-flow control, and capability-based security to accommodate LLM non-determinism is a key open problem. Progent introduces the first privilege control framework for LLM agents, reducing ASR to 0\% on benchmarks~\cite{shi2025progent}. CELLMATE bridges the semantic gap between browser actions and privilege requirements~\cite{meng2025cellmate}. Mandatory access control assumes deterministic behavior incompatible with LLMs; no consensus exists on capability granularity for AI agents.
\end{enumerate}


%% ------------------------------------------------------------
\subsection{Unified Multi-Modal Detection Frameworks}
\label{sec:future-detection}
%% ------------------------------------------------------------

\emph{Gap:} Most detectors operate on a single modality independently. No single framework detects AI-generated content across text, image, audio, and video within a unified architecture. The closest attempt, BusterX++, covers only images and video~\cite{wen2025busterx}. Audio-visual detectors (AVFF~\cite{oorloff2024avff}) and image-text manipulation detectors (DGM$^4$~\cite{shao2023dgm4}) address paired modalities but not the full spectrum.

\emph{Why it matters:} Real-world disinformation combines multiple modalities. Single-modality detectors miss cross-modal inconsistencies and can be circumvented by manipulating unmonitored modalities. State-of-the-art detectors suffer AUC drops of 45--50\% on in-the-wild content versus academic benchmarks~\cite{chandra2025deepfakeeval}.

\emph{Research directions:} Foundation models show promise as detection backbones---CLIP's feature space encodes transferable ``realness'' properties~\cite{ojha2023univfd}. Cross-modal transfer learning for detection remains entirely unexplored: no study has identified universal artifacts shared across modalities. Research paths include developing modality-agnostic generation fingerprints, building unified foundation models pre-trained on multi-modal forensic tasks, and establishing shared artifact representations grounded in the common mathematical structure of diffusion and autoregressive generation.


%% ------------------------------------------------------------
\subsection{Provably Robust Watermarking}
\label{sec:future-watermark}
%% ------------------------------------------------------------

\emph{Gap:} All practical watermarking schemes rely on empirical robustness evaluation. No deployed scheme has formal security proofs against adaptive adversaries, and a fundamental impossibility result constrains what is achievable.

Zhang et al.\ prove that under natural assumptions---access to a quality oracle and a perturbation oracle---\textbf{strong watermarking is fundamentally impossible}: a generic quality-preserving random-walk attack removes watermarks from KGW and related schemes~\cite{zhang2024impossibility}. Pang et al.\ identify a critical robustness-spoofing tradeoff: robust watermarks are inherently vulnerable to spoofing attacks~\cite{pang2024nofree}. Jovanovi\'c et al.\ demonstrate watermark stealing via API queries~\cite{jovanovic2024stealing}.

On the constructive side, Christ et al.\ achieve \emph{undetectable} watermarks based on one-way functions~\cite{christ2024undetectable}. Golowich and Moitra achieve both undetectability and edit-distance robustness~\cite{golowich2024editdistance}. Zhao et al.\ provide the first provably robust guarantees with exponentially decaying error rates~\cite{zhao2024provable}. Google's SynthID has been validated on $\sim$20 million chatbot responses~\cite{dathathri2024synthid}.

\emph{What remains unsolved:} The impossibility result means all practical watermarks can theoretically be removed. Low-entropy text (factual answers, code) is fundamentally unwatermarkable. No cross-provider interoperable standard exists. Research paths include tightening information-theoretic bounds for realistic token generation, developing efficient cryptographic watermark instantiations, exploring semantic watermarks robust to paraphrase attacks, and establishing cross-provider standards complementary to C2PA.


%% ------------------------------------------------------------
\subsection{Privacy-Preserving Generation}
\label{sec:future-privacy}
%% ------------------------------------------------------------

\emph{Gap:} Generative models memorize and can regurgitate training data, and all known defenses impose severe utility costs.

Nasr et al.\ demonstrate that gigabytes of training data can be extracted from production models including ChatGPT, using a divergence attack causing emission at 150$\times$ the normal rate~\cite{nasr2025extraction}. Memorization grows log-linearly with model capacity, data duplication, and prompt context length~\cite{carlini2023quantifying}. For diffusion models, over 1,000 training images were extracted from Stable Diffusion~\cite{carlini2023extractingdiffusion}. Critically, blocking verbatim memorization provides a false sense of privacy---style-transfer prompts easily circumvent filters~\cite{ippolito2023verbatim}.

Differential privacy faces fundamental scalability challenges. Tram\`er et al.'s ICML 2024 Best Paper critically challenges the public-pretraining-then-private-fine-tuning paradigm: web-scraped ``public'' data contains sensitive information, and DP pretraining from scratch at frontier scale remains computationally infeasible~\cite{tramer2024dpposition}.

Machine unlearning offers no reliable solution. The TOFU benchmark demonstrates that no tested unlearning baseline makes models behave as if target data was never seen~\cite{maini2024tofu}. Even exact unlearning---retraining from scratch---still leaks sensitive information when adversaries access pre-unlearning checkpoints~\cite{unlearned2025}.

\emph{Research paths:} Memorization-aware training procedures without full DP overhead; tight privacy auditing tools closing the gap between theoretical and empirical $\varepsilon$; unlearning verification with adversarial guarantees; inference-time privacy mechanisms offering new utility-privacy frontiers.


%% ------------------------------------------------------------
\subsection{International Governance Harmonization}
\label{sec:future-governance}
%% ------------------------------------------------------------

\emph{Gap:} Three major regulatory regimes---the EU AI Act, China's generative AI regulations, and US executive actions---have adopted fundamentally different approaches, creating jurisdictional conflicts, regulatory arbitrage, and standards fragmentation with no operational mutual recognition mechanism.

The EU mandates machine-readable watermarking by August 2026~\cite{euaiact2024}; China was the first to enact binding generative AI regulation~\cite{chinainterim2023}, requiring alignment with ``Core Socialist Values'' and mandatory AI content labeling; the US revoked its primary AI safety executive order within hours of a new administration~\cite{eo14179}. The Council of Europe adopted the first legally binding international AI treaty (CETS No.~225, May 2024), signed by the US, EU, and UK but not China~\cite{coe2024treaty}. The G7 Hiroshima Process produced international guiding principles~\cite{hiroshima2023}, and the UN established an Independent International Scientific Panel on AI~\cite{unga2025aipanel}.

As UNU Rector Marwala warns: ``We are entering an era of AI governance arbitrage, in which the pace of innovation hinges more on how well actors navigate regulatory gaps than on technological breakthroughs''~\cite{marwala2026arbitrage}.

\emph{What remains unsolved:} No mutual recognition mechanism for AI safety certifications is operational. Content labeling standards are fragmented across jurisdictions. National security carve-outs create blind spots. Research paths include interoperable technical standards for content provenance, mutual recognition frameworks analogous to pharmaceutical regulatory cooperation, cross-border regulatory sandboxes, and empirical metrics for comparing regulatory effectiveness.


%% ------------------------------------------------------------
\subsection{Dynamic Evaluation Frameworks}
\label{sec:future-eval}
%% ------------------------------------------------------------

\emph{Gap:} Static benchmarks become obsolete through saturation and contamination. LLMs now exceed 90\% on benchmarks like MMLU and HellaSwag, and benchmark data increasingly appears in training sets. Evaluation methods for agentic and multi-modal systems remain immature.

\emph{Living benchmarks.} LiveBench~\cite{white2024livebench} introduces monthly-updated questions where top models achieve under 70\% accuracy. Humanity's Last Exam contributes 2,500 expert-level questions with a dynamic rolling fork~\cite{phan2025humanitylastexam}. However, Sun et al.\ demonstrate a fundamental fidelity-resistance tradeoff: surface-level edits fail to prevent memorization, while deeper rewrites distort evaluation tasks~\cite{sun2025contamination}.

\emph{Agentic evaluation.} AgentBench evaluates across 8 real-world environments~\cite{liu2024agentbench}. Agent-SafetyBench reveals that no tested agent achieves a safety score above 60\% across 349 interaction environments~\cite{zhang2024agentsafetybench}. Agent Security Bench formalizes attacks and defenses across 10 agentic scenarios~\cite{asb2025}.

\emph{Community safety infrastructure.} MLCommons AILuminate v1.0 provides the first industry-standard safety benchmark with 24,000+ prompts across 12 hazard categories~\cite{ghosh2024ailuminate}. HELM offers a foundational living benchmark framework with 7 metrics across 16 scenarios~\cite{liang2023helm}. WildTeaming mines 5,700 unique jailbreak tactic clusters from in-the-wild interactions~\cite{jiang2024wildteaming}.

\emph{What remains unsolved:} Contamination-proof evaluation remains an open problem. Unified cross-modal safety evaluation is nascent. Long-horizon agentic safety evaluation lacks principled methods. Most benchmarks are English-centric. Most fundamentally, benchmark performance does not reliably predict real-world safety---MLCommons explicitly notes ``negative predictive power'' where passing benchmarks does not ensure deployment safety. Research paths include contamination-resistant evaluation via procedural generation, agentic evaluation with verified sandboxes, culturally diverse adversarial corpora, and formal relationships between benchmark and deployment-time safety.
